\section{Introduction}
\subsection{Problem background}
Selecting a material for solar water splitting requires several properties to agree: favorable formation energy, an appropriate band gap, and suitable band-edge positions. Computational screening of cubic perovskites provides a composition space in which these requirements can be studied together~\cite{castelli2012}. In the supplied problem, five penalties combine their departures from the target conditions into a violation degree. Predicting this quantity from composition allows unlabeled candidates to be compared and provides the objective values needed for subsequent discrete selection.

\subsection{Problem restatement}
\begin{itemize}
  \item \textbf{Problem 1: Surrogate prediction.} Construct a composition-based predictor of the violation degree using the supplied training data. Assess numerical accuracy and the ability to prioritize low-violation materials, select the model using the validation data, and predict the unlabeled test compositions.
  \item \textbf{Problem 2:}
  \item \textbf{Problem 3:}
\end{itemize}

\section{Problem Analysis}
\subsection{Analysis of Problem 1}
\begin{itemize}
  \item \textbf{Imbalanced physical regimes.} Zero band gaps identify metals and must be retained. Their unavoidable band-gap penalty motivates separate conditional regressions within one composition-based predictor.
  \item \textbf{Site-specific chemistry.} A and B have distinct roles, while X denotes a three-atom group. Site descriptors, X stoichiometry and their contrasts are needed to distinguish these compositions.
  \item \textbf{Two prediction requirements.} Low average error does not by itself establish useful screening. Numerical error and low-violation enrichment therefore require separate measures under validation that keeps each ordered A--B pair together.
\end{itemize}
\subsection{Analysis of Problem 2}
\subsection{Analysis of Problem 3}

\section{Model Assumptions and Notation}
\subsection{Core assumptions}
\begin{enumerate}
  \item \textbf{A common material setting.} Composition and fixed elemental descriptors are assumed sufficient for prediction within the supplied cubic setting. This does not cover other structures or experimental synthesis conditions.
  \item \textbf{Transfer within represented chemistry.} Relations learned from training compositions are assumed useful for new A--B pairs in the supplied domain. Grouped validation and exclusion tests examine this assumption.
\end{enumerate}

\subsection{Principal notation}
Table~\ref{tab:notation} collects the principal symbols. Each is also defined when first used in the model.
\par
\begin{table}[H]
\centering
\caption{Principal notation}\label{tab:notation}
\begin{tabularx}{\linewidth}{@{}l X l@{}}
\toprule
Symbol & Meaning & Unit or range\\
\midrule
$c=(a,b,x)$ & Ordered composition at sites A, B and group X & Supplied grid\\
$\boldsymbol\phi(c)$ & Composition-derived descriptor vector & 336 entries\\
$V,\widehat V$ & Observed and predicted violation degree & Dimensionless\\
$G$ & Nonmetal indicator, $\mathbf{1}\{E_g>0\}$ & $0$ or $1$\\
$\widehat p$ & Estimated conditional probability of $G=1$ & $[0,1]$\\
$\widehat\mu_0,\widehat\mu_1$ & Metal and nonmetal conditional predictions & Dimensionless\\
$\varepsilon,k$ & Near-positive cutoff and selection budget & $2$, $20$\\
$N,M,h_k$ & Pool size, near-positive count and selected hits & Counts\\
\bottomrule
\end{tabularx}
\end{table}
\par
\FloatBarrier
